\documentclass[11pt]{article}
\usepackage{polynom}
\usepackage[a4paper,margin=0.75in]{geometry}
\usepackage{fancyhdr}
\usepackage{enumitem}
\usepackage{amsmath}
\usepackage{amssymb}
\usepackage{graphicx}
\usepackage{natbib}

\usepackage{tikz}
\usepackage{ytableau}

\title{{\textbf{STA380 Project Proposal}}

\vspace{0.5 cm}

\Large Computing Empirical and Theoretical Autocorrelation of Time Series Models}
\author{\textbf{Daniel Cowley, Arnav Nagzirkar, Esteban Gonzalez, } \\
        \textbf{Devesh Aggarwal, Khushal Patel, Sohaib Sajid}}
\date{}

\begin{document}

\maketitle

\section{Project Topic}

We will be computing the sample and theoretical autocovariance and autocorrelation for the following time series models: 
\begin{itemize} 
    \item Autoregressive (AR) 
    \item Moving Average (MA) 
    \item Autoregressive Moving Average (ARMA) 
\end{itemize}

\section{Simulation vs. Dataset}

The results will only be from simulations.

\section{Project Details}

This project will focus on generating data from AR($k$), MA($k$), and ARMA($p,q$) time series models. We will compute both empirical and theoretical autocovariance and autocorrelation
functions (through our own implementations), and compare the results. The time series processes will be generated from either user input or random coefficients.

\vspace{0.5 cm}

\noindent Specification of models and their parameters: (\cite{johnson2009})
\begin{itemize}
    \item Autoregressive (AR($k$))
    \[Y_t = \sum_{i = 1}^k a_iY_{i - 1} + w_t, \quad w_t \overset{\text{i.i.d.}}{\sim} N(0, \sigma_W^2), \quad k \in \{1, 2, 3\}\]

    Note: Theoretical values for autocorrelation and autocovariance can be generated recursively using the Yule-Walker equations.

    \item Moving Average (MA($k$))
    \[Y_t = w_t + \sum_{i = 1}^k b_iw_{t - i}, \quad w_t \overset{\text{i.i.d.}}{\sim} N(0, \sigma_W^2), \quad k \in \{1, 2, 3\}\]

    \item Autoregressive Moving Average (ARMA($p, q$))
    \[Y_t = \sum_{i = 1}^p a_iY_{i - 1} + w_t + \sum_{i = 1}^q b_iw_{t - i}, \quad w_t \overset{\text{i.i.d.}}{\sim} N(0, \sigma_W^2), \quad p \in \{1, 2, 3\}, \quad q \in \{1, 2, 3\}\]
\end{itemize}

\newpage

\noindent The following outputs will be provided:
\begin{itemize}
    \item Plots of the generated time series
    \item Sample ACVF and ACF plots alongside their theoretical counterparts
    \item Bias between the sampled and true ACVF and ACF plots
\end{itemize}

\section{User Inputs (Shiny Components)}

Users will be able to modify the following:
\begin{enumerate}
    \item Select which time series model to simulate.
    
    \item Toggle the display of the time series, ACF, and ACVF plots (in any configuration from individually to all together).

    \item Ability to overlay the theoretical and sampled counterparts on the ACF and ACVF plots.

    \item Specify model coefficients or generate them randomly.

    \item Specify values of $k$ (up to at least 5).

    \item Set the simulation length $t$ (to adjust computational time and accuracy of simulated stationarity conditions).

    \item Specify a seed

    \item Customize the color scheme of various UI elements (time permitted).

    \item Ability to compare the computational efficiency of our methods with existing implementations (time permitted).
\end{enumerate}

\nocite{tidyverse}
\nocite{shiny}
\nocite{tseries}
\nocite{testthat}

\bibliographystyle{apalike}
\bibliography{references}

\end{document}
